\documentclass{article}
\usepackage[preprint]{neurips_2026}

\usepackage[utf8]{inputenc}
\usepackage[T1]{fontenc}
\usepackage{hyperref}
\usepackage{url}
\usepackage{booktabs}
\usepackage{amsfonts}
\usepackage{nicefrac}
\usepackage{microtype}
\usepackage{xcolor}
\usepackage{graphicx}
\usepackage{amsmath}
\usepackage{array}
\usepackage{longtable}
\usepackage{caption}

\title{基于 3DGS 与 AIGC 的多源资产生成与真实场景融合}

\author{%
  姓名待填写 \\
  学号待填写 \\
  课程：深度学习与空间智能 \\
  \texttt{email@example.com}
}

\begin{document}

\maketitle

\begin{abstract}
本文围绕课程作业题目一，构建一个从真实场景重建、AIGC 资产生成到统一三维融合渲染的完整实验流程。我们将分别完成基于多视角真实素材的物体 A 重建、基于文本提示的物体 B 生成、基于单张图像的物体 C 生成，并选取开源真实场景作为统一背景。随后，我们将三类前景资产统一转换为高斯表示，与背景 3DGS 场景进行拼接和漫游渲染。本文给出工程结构、关键配置、训练与导出流程，并为后续补充实验结果、指标曲线和失败分析预留统一入口。
\end{abstract}

\section{任务定义}
本项目严格对应课程作业题目一，要求同时完成以下四类结果：
\begin{enumerate}
  \item 物体 A：真实多视角素材，经 COLMAP 位姿恢复后使用 3D Gaussian Splatting 进行重建；
  \item 物体 B：仅由文本提示驱动的文本到 3D 生成；
  \item 物体 C：由单张真实图像驱动的单图到 3D 生成；
  \item 独立背景场景：从公开数据集选取一个真实场景并使用 3DGS 重建。
\end{enumerate}
在此基础上，三类前景资产需要与背景统一融合，输出多视角漫游视频，并在报告中给出方法、配置、结果与分析。

\section{方法总览}
\subsection{整体流程}
整体工程分为 Windows 本地准备与 Linux 集群执行两个阶段。Windows 侧负责准备源码、离线权重、输入素材、交接文档与报告骨架；Linux 侧负责正式训练、结果导出和融合渲染。

\subsection{物体 A：真实多视角重建}
物体 A 采用公开真实多视角图像或后续替换为自采素材，先使用 COLMAP 执行特征提取、匹配与稀疏重建，再将结果交给 3DGS 优化器进行训练。该路线严格满足题意中“真实多视角 + COLMAP + 3DGS”的要求。

\subsection{物体 B：文本到 3D}
物体 B 使用 threestudio 主线实现文本到 3D。默认提示词为一个几何边界清晰、材质细节适中的静态物体，以提高训练稳定性。如果出现依赖阻塞，仅允许在 threestudio 内切换为同类文本到 3D 配置，不允许退化为占位式 mesh。

\subsection{物体 C：单图到 3D}
物体 C 使用 Stable Zero123 主线，输入单张前景图像及其去背景版本，输出 3D 资产。若实现细节需要调整，也只能在单图到 3D 体系内部选择兼容路线。

\subsection{统一高斯融合}
为了保持最终表示一致性，本项目不采用“背景是高斯、前景直接是网格”的最终拼接方式，而是将 B、C 的输出进一步转换为可被 3DGS 重建的统一视图数据，再分别重建为前景高斯场。最终，A、B、C 与背景都以高斯表示存在，并通过刚体变换、尺度匹配和位置布局完成统一融合。

\section{实验设置}
\subsection{硬件与运行环境}
默认执行环境为 Linux 集群单卡 NVIDIA A100 40GB。由于集群网络能力有限，本项目采用“本地预打包不可达资产 + 集群在线安装可达依赖”的混合部署方式。

\subsection{默认数据与输入}
\begin{itemize}
  \item 背景场景：Mip-NeRF 360 中的 garden，备选为 counter。
  \item 物体 A：公开真实多视角素材，后续可无缝替换为自采素材。
  \item 物体 B：预定义文本提示词集合。
  \item 物体 C：单张主体清晰的前景图像及其 RGBA 版本。
\end{itemize}

\subsection{默认日志与指标}
由于 WandB 不可用，训练日志统一落地为本地文件，包括：
\begin{itemize}
  \item \texttt{train.log}
  \item \texttt{metrics.json}
  \item \texttt{curves.csv}
  \item \texttt{gpu\_mem.txt}
\end{itemize}
最终报告将从这些文件中提取损失曲线、重建指标、训练时间和显存占用。

\section{工程结构与执行流程}
\subsection{脚本入口}
仓库提供如下稳定入口：
\begin{itemize}
  \item \texttt{scripts/setup\_cluster.sh}
  \item \texttt{scripts/prepare\_assets.py}
  \item \texttt{scripts/train\_bg\_3dgs.sh}
  \item \texttt{scripts/train\_objA\_3dgs.sh}
  \item \texttt{scripts/train\_objB\_text\_to\_3d.sh}
  \item \texttt{scripts/train\_objC\_image\_to\_3d.sh}
  \item \texttt{scripts/convert\_to\_gaussians.py}
  \item \texttt{scripts/merge\_gaussians.py}
  \item \texttt{scripts/render\_flythrough.sh}
  \item \texttt{scripts/export\_report\_assets.py}
\end{itemize}

\subsection{执行顺序}
推荐执行顺序为：环境准备、资产校验、背景训练、物体 A、物体 B、物体 C、统一高斯化、融合渲染、报告资产导出。所有阶段都支持独立重跑，不要求捆绑成单次大脚本。

\section{结果展示占位}
\begin{figure}[t]
  \centering
  \fbox{\rule{0pt}{1.8in}\rule{0.85\linewidth}{0pt}}
  \caption{最终融合场景主视图占位图。后续替换为融合后渲染结果。}
  \label{fig:fusion_teaser}
\end{figure}

\begin{table}[t]
  \centering
  \caption{实验配置与关键指标占位表。}
  \begin{tabular}{lccc}
    \toprule
    模块 & 关键配置 & 主要指标 & 备注 \\
    \midrule
    背景 3DGS & 待填 & 待填 & 待填 \\
    物体 A & 待填 & 待填 & 待填 \\
    物体 B & 待填 & 待填 & 待填 \\
    物体 C & 待填 & 待填 & 待填 \\
    \bottomrule
  \end{tabular}
  \label{tab:summary}
\end{table}

\section{讨论与风险}
本项目的主要风险来自三方面：其一，文本到 3D 与单图到 3D 依赖链较长，对权重路径和 CUDA 扩展敏感；其二，集群对部分外部站点不可达，因此必须提前准备离线资产；其三，最终统一高斯化与场景融合需要额外的中间转换和对齐步骤。为保证题意严格性，本项目只允许在同类方法内部切换 fallback，不允许退化为占位式伪实现。

\section{结论与后续工作}
当前版本首先完成了从零搭建的可交接工程骨架，为后续在 Linux 集群上执行正式实验提供了统一结构、脚本接口和报告模板。后续工作将集中于补齐离线资产、执行完整训练、生成统一融合场景，并回填本报告中的结果图、定量指标与失败案例分析。

\bibliographystyle{plainnat}
\bibliography{references}

\end{document}
